\IfFileExists{stacks-project.cls}{%
\documentclass{stacks-project}
}{%
\documentclass{amsart}
}

% For dealing with references we use the comment environment
\usepackage{verbatim}
\newenvironment{reference}{\comment}{\endcomment}
%\newenvironment{reference}{}{}
\newenvironment{slogan}{\comment}{\endcomment}
\newenvironment{history}{\comment}{\endcomment}

% For commutative diagrams we use Xy-pic
\usepackage[all]{xy}

% We use 2cell for 2-commutative diagrams.
\xyoption{2cell}
\UseAllTwocells

% We use multicol for the list of chapters between chapters
\usepackage{multicol}

% This is generally recommended for better output
\usepackage{lmodern}
\usepackage[T1]{fontenc}

% For cross-file-references
\usepackage{xr-hyper}

% Package for hypertext links:
\usepackage{hyperref}

% For any local file, say "hello.tex" you want to link to please
% use \externaldocument[hello-]{hello}
\externaldocument[introduction-]{introduction}
\externaldocument[conventions-]{conventions}
\externaldocument[sets-]{sets}
\externaldocument[categories-]{categories}
\externaldocument[topology-]{topology}
\externaldocument[sheaves-]{sheaves}
\externaldocument[sites-]{sites}
\externaldocument[stacks-]{stacks}
\externaldocument[fields-]{fields}
\externaldocument[algebra-]{algebra}
\externaldocument[brauer-]{brauer}
\externaldocument[homology-]{homology}
\externaldocument[derived-]{derived}
\externaldocument[simplicial-]{simplicial}
\externaldocument[more-algebra-]{more-algebra}
\externaldocument[smoothing-]{smoothing}
\externaldocument[modules-]{modules}
\externaldocument[sites-modules-]{sites-modules}
\externaldocument[injectives-]{injectives}
\externaldocument[cohomology-]{cohomology}
\externaldocument[sites-cohomology-]{sites-cohomology}
\externaldocument[dga-]{dga}
\externaldocument[dpa-]{dpa}
\externaldocument[sdga-]{sdga}
\externaldocument[hypercovering-]{hypercovering}
\externaldocument[schemes-]{schemes}
\externaldocument[constructions-]{constructions}
\externaldocument[properties-]{properties}
\externaldocument[morphisms-]{morphisms}
\externaldocument[coherent-]{coherent}
\externaldocument[divisors-]{divisors}
\externaldocument[limits-]{limits}
\externaldocument[varieties-]{varieties}
\externaldocument[topologies-]{topologies}
\externaldocument[descent-]{descent}
\externaldocument[perfect-]{perfect}
\externaldocument[more-morphisms-]{more-morphisms}
\externaldocument[flat-]{flat}
\externaldocument[groupoids-]{groupoids}
\externaldocument[more-groupoids-]{more-groupoids}
\externaldocument[etale-]{etale}
\externaldocument[chow-]{chow}
\externaldocument[intersection-]{intersection}
\externaldocument[pic-]{pic}
\externaldocument[weil-]{weil}
\externaldocument[adequate-]{adequate}
\externaldocument[dualizing-]{dualizing}
\externaldocument[duality-]{duality}
\externaldocument[discriminant-]{discriminant}
\externaldocument[derham-]{derham}
\externaldocument[local-cohomology-]{local-cohomology}
\externaldocument[algebraization-]{algebraization}
\externaldocument[curves-]{curves}
\externaldocument[resolve-]{resolve}
\externaldocument[models-]{models}
\externaldocument[functors-]{functors}
\externaldocument[equiv-]{equiv}
\externaldocument[pione-]{pione}
\externaldocument[etale-cohomology-]{etale-cohomology}
\externaldocument[proetale-]{proetale}
\externaldocument[relative-cycles-]{relative-cycles}
\externaldocument[more-etale-]{more-etale}
\externaldocument[trace-]{trace}
\externaldocument[crystalline-]{crystalline}
\externaldocument[spaces-]{spaces}
\externaldocument[spaces-properties-]{spaces-properties}
\externaldocument[spaces-morphisms-]{spaces-morphisms}
\externaldocument[decent-spaces-]{decent-spaces}
\externaldocument[spaces-cohomology-]{spaces-cohomology}
\externaldocument[spaces-limits-]{spaces-limits}
\externaldocument[spaces-divisors-]{spaces-divisors}
\externaldocument[spaces-over-fields-]{spaces-over-fields}
\externaldocument[spaces-topologies-]{spaces-topologies}
\externaldocument[spaces-descent-]{spaces-descent}
\externaldocument[spaces-perfect-]{spaces-perfect}
\externaldocument[spaces-more-morphisms-]{spaces-more-morphisms}
\externaldocument[spaces-flat-]{spaces-flat}
\externaldocument[spaces-groupoids-]{spaces-groupoids}
\externaldocument[spaces-more-groupoids-]{spaces-more-groupoids}
\externaldocument[bootstrap-]{bootstrap}
\externaldocument[spaces-pushouts-]{spaces-pushouts}
\externaldocument[spaces-chow-]{spaces-chow}
\externaldocument[groupoids-quotients-]{groupoids-quotients}
\externaldocument[spaces-more-cohomology-]{spaces-more-cohomology}
\externaldocument[spaces-simplicial-]{spaces-simplicial}
\externaldocument[spaces-duality-]{spaces-duality}
\externaldocument[formal-spaces-]{formal-spaces}
\externaldocument[restricted-]{restricted}
\externaldocument[spaces-resolve-]{spaces-resolve}
\externaldocument[formal-defos-]{formal-defos}
\externaldocument[defos-]{defos}
\externaldocument[cotangent-]{cotangent}
\externaldocument[examples-defos-]{examples-defos}
\externaldocument[algebraic-]{algebraic}
\externaldocument[examples-stacks-]{examples-stacks}
\externaldocument[stacks-sheaves-]{stacks-sheaves}
\externaldocument[criteria-]{criteria}
\externaldocument[artin-]{artin}
\externaldocument[quot-]{quot}
\externaldocument[stacks-properties-]{stacks-properties}
\externaldocument[stacks-morphisms-]{stacks-morphisms}
\externaldocument[stacks-limits-]{stacks-limits}
\externaldocument[stacks-cohomology-]{stacks-cohomology}
\externaldocument[stacks-perfect-]{stacks-perfect}
\externaldocument[stacks-introduction-]{stacks-introduction}
\externaldocument[stacks-more-morphisms-]{stacks-more-morphisms}
\externaldocument[stacks-geometry-]{stacks-geometry}
\externaldocument[moduli-]{moduli}
\externaldocument[moduli-curves-]{moduli-curves}
\externaldocument[examples-]{examples}
\externaldocument[exercises-]{exercises}
\externaldocument[guide-]{guide}
\externaldocument[desirables-]{desirables}
\externaldocument[coding-]{coding}
\externaldocument[obsolete-]{obsolete}
\externaldocument[fdl-]{fdl}
\externaldocument[index-]{index}

% Theorem environments.
%
\theoremstyle{plain}
\newtheorem{theorem}[subsection]{Theorem}
\newtheorem{proposition}[subsection]{Proposition}
\newtheorem{lemma}[subsection]{Lemma}

\theoremstyle{definition}
\newtheorem{definition}[subsection]{Definition}
\newtheorem{example}[subsection]{Example}
\newtheorem{exercise}[subsection]{Exercise}
\newtheorem{situation}[subsection]{Situation}

\theoremstyle{remark}
\newtheorem{remark}[subsection]{Remark}
\newtheorem{remarks}[subsection]{Remarks}

\numberwithin{equation}{subsection}

% Macros
%
\def\lim{\mathop{\mathrm{lim}}\nolimits}
\def\colim{\mathop{\mathrm{colim}}\nolimits}
\def\Spec{\mathop{\mathrm{Spec}}}
\def\Hom{\mathop{\mathrm{Hom}}\nolimits}
\def\Ext{\mathop{\mathrm{Ext}}\nolimits}
\def\SheafHom{\mathop{\mathcal{H}\!\mathit{om}}\nolimits}
\def\SheafExt{\mathop{\mathcal{E}\!\mathit{xt}}\nolimits}
\def\Sch{\mathit{Sch}}
\def\Mor{\mathop{\mathrm{Mor}}\nolimits}
\def\Ob{\mathop{\mathrm{Ob}}\nolimits}
\def\Sh{\mathop{\mathit{Sh}}\nolimits}
\def\NL{\mathop{N\!L}\nolimits}
\def\CH{\mathop{\mathrm{CH}}\nolimits}
\def\proetale{{pro\text{-}\acute{e}tale}}
\def\etale{{\acute{e}tale}}
\def\QCoh{\mathit{QCoh}}
\def\Ker{\mathop{\mathrm{Ker}}}
\def\Im{\mathop{\mathrm{Im}}}
\def\Coker{\mathop{\mathrm{Coker}}}
\def\Coim{\mathop{\mathrm{Coim}}}

% Boxtimes
%
\DeclareMathSymbol{\boxtimes}{\mathbin}{AMSa}{"02}

%
% Macros for moduli stacks/spaces
%
\def\QCohstack{\mathcal{QC}\!\mathit{oh}}
\def\Cohstack{\mathcal{C}\!\mathit{oh}}
\def\Spacesstack{\mathcal{S}\!\mathit{paces}}
\def\Quotfunctor{\mathrm{Quot}}
\def\Hilbfunctor{\mathrm{Hilb}}
\def\Curvesstack{\mathcal{C}\!\mathit{urves}}
\def\Polarizedstack{\mathcal{P}\!\mathit{olarized}}
\def\Complexesstack{\mathcal{C}\!\mathit{omplexes}}
% \Pic is the operator that assigns to X its picard group, usage \Pic(X)
% \Picardstack_{X/B} denotes the Picard stack of X over B
% \Picardfunctor_{X/B} denotes the Picard functor of X over B
\def\Pic{\mathop{\mathrm{Pic}}\nolimits}
\def\Picardstack{\mathcal{P}\!\mathit{ic}}
\def\Picardfunctor{\mathrm{Pic}}
\def\Deformationcategory{\mathcal{D}\!\mathit{ef}}
\usepackage{cleveref}
\crefname{originalreference}{original reference}{original references}
\Crefname{originalreference}{Original reference}{Original references}
\makeatletter
\newlabel{categories-definition-multiplicative-system}{{04VC}{source}{}{sourcepointer-1}{}}
\newlabel{categories-definition-multiplicative-system@cref}{{[originalreference][0][]04VC}{[1][0][]source}{}{}{}}
\makeatother
\crefname{originalreference}{original reference}{original references}
\Crefname{originalreference}{Original reference}{Original references}
\makeatletter
\newlabel{categories-lemma-almost-finite-limits-exist}{{04AU}{source}{}{sourcepointer-2}{}}
\newlabel{categories-lemma-almost-finite-limits-exist@cref}{{[originalreference][0][]04AU}{[1][0][]source}{}{}{}}
\makeatother
\crefname{originalreference}{original reference}{original references}
\Crefname{originalreference}{Original reference}{Original references}
\makeatletter
\newlabel{categories-lemma-composition-cartesian}{{02XL}{source}{}{sourcepointer-3}{}}
\newlabel{categories-lemma-composition-cartesian@cref}{{[originalreference][0][]02XL}{[1][0][]source}{}{}{}}
\makeatother
\title{Strongly cartesian identity-base morphisms form a right multiplicative system}
\author{The Stacks Project Authors}
\date{stacks-project:stacks.tex:local-snapshot-20261012; GNU FDL 1.2 or later}
\begin{document}
\maketitle
\noindent\textit{Editorial note.} No omission markers in-span; proof complete as written.\par\medskip
\noindent\textit{Allowed standard prerequisites (editorial).}\begin{itemize}
\item Right multiplicative system axioms. categories-definition-multiplicative-system (cross-chapter, named). Fixes RMS1, RMS2, RMS3 to be verified.
\item Composition of strongly cartesian morphisms. categories-lemma-composition-cartesian (cross-chapter, named). Settles RMS1 (identity is in R, compositions stay in R).
\item Equalizers in the base. categories-lemma-almost-finite-limits-exist (cross-chapter, named). Provides the equalizer of a, a' used to build the RMS3 dominating arrow.
\end{itemize}\medskip
\ifdefined\hypertarget\hypertarget{sourcepointer-1}{}\fi
\noindent\textit{Source reference.} Original reference 04VC: focused-seeds/geometry/S9-stk-right-multiplicative-system/ref-tags.txt, source lines 5--5. Navigation pointer only: the counted proof cites this statement by \textbackslash{}ref as a standing input; the deduction the counted proof has to make is carried out inside the counted span, and no part of that deduction is delegated to this pointer.\par\medskip
\ifdefined\hypertarget\hypertarget{sourcepointer-2}{}\fi
\noindent\textit{Source reference.} Original reference 04AU: focused-seeds/geometry/S9-stk-right-multiplicative-system/ref-tags.txt, source lines 6--6. Navigation pointer only: the counted proof cites this statement by \textbackslash{}ref as a standing input; the deduction the counted proof has to make is carried out inside the counted span, and no part of that deduction is delegated to this pointer.\par\medskip
\ifdefined\hypertarget\hypertarget{sourcepointer-3}{}\fi
\noindent\textit{Source reference.} Original reference 02XL: focused-seeds/geometry/S9-stk-right-multiplicative-system/ref-tags.txt, source lines 7--7. Navigation pointer only: the counted proof cites this statement by \textbackslash{}ref as a standing input; the deduction the counted proof has to make is carried out inside the counted span, and no part of that deduction is delegated to this pointer.\par\medskip
\begin{lemma}
\label{lemma-right-multiplicative-system}
In the situation above assume
\begin{enumerate}
\item $p : \mathcal{S} \to \mathcal{C}$ is a fibred category,
\item $\mathcal{C}$ has nonempty finite limits, and
\item $u : \mathcal{C} \to \mathcal{D}$ commutes with nonempty finite limits.
\end{enumerate}
Consider the set $R \subset \text{Arrows}(u_{pp}\mathcal{S})$ of morphisms
of the form
$$
(a, \text{id}_V, \alpha) :
(U', \phi' : V \to u(U'), x')
\longrightarrow
(U, \phi : V \to u(U), x)
$$
with $\alpha$ strongly cartesian. Then $R$ is a right multiplicative system.
\end{lemma}

\begin{proof}
According to
Categories, Definition \ref{categories-definition-multiplicative-system}
we have to check RMS1, RMS2, RMS3.
Condition RMS1 holds as a composition of strongly cartesian morphisms is
strongly cartesian, see
Categories, Lemma \ref{categories-lemma-composition-cartesian}.

\medskip\noindent
To check RMS2 suppose we have a morphism
$$
(a, b, \alpha) :
(U_1, \phi_1 : V_1 \to u(U_1), x_1)
\longrightarrow
(U, \phi : V \to u(U), x)
$$
of $u_{pp}\mathcal{S}$ and a morphism
$$
(c, \text{id}_V, \gamma) :
(U', \phi' : V \to u(U'), x')
\longrightarrow
(U, \phi : V \to u(U), x)
$$
with $\gamma$ strongly cartesian from $R$. In this situation set
$U'_1 = U_1 \times_U U'$, and denote $a' : U'_1 \to U'$ and
$c' : U'_1 \to U_1$ the projections.
As $u(U'_1) = u(U_1) \times_{u(U)} u(U')$
we see that $\phi'_1 = (\phi_1, \phi') : V_1 \to u(U'_1)$ is
a morphism in $\mathcal{D}$. Let $\gamma_1 : x_1' \to x_1$
be a strongly cartesian morphism of $\mathcal{S}$ with
$p(\gamma_1) = \phi'_1$ (which exists because $\mathcal{S}$ is a
fibred category over $\mathcal{C}$). Then as $\gamma : x' \to x$ is
strongly cartesian there exists a unique morphism
$\alpha' : x'_1 \to x'$ with $p(\alpha') = a'$.
At this point we see that
$$
(a', b, \alpha') :
(U_1, \phi_1 : V_1 \to u(U'_1), x'_1)
\longrightarrow
(U, \phi : V \to u(U'), x')
$$
is a morphism and that
$$
(c', \text{id}_{V_1}, \gamma_1) :
(U'_1, \phi'_1 : V_1 \to u(U'_1), x'_1)
\longrightarrow
(U_1, \phi : V_1 \to u(U_1), x_1)
$$
is an element of $R$ which form a solution of the existence problem
posed by RMS2.

\medskip\noindent
Finally, suppose that
$$
(a, b, \alpha), (a', b', \alpha') :
(U_1, \phi_1 : V_1 \to u(U_1), x_1)
\longrightarrow
(U, \phi : V \to u(U), x)
$$
are two morphisms of $u_{pp}\mathcal{S}$ and suppose that
$$
(c, \text{id}_V, \gamma) :
(U, \phi : V \to u(U), x)
\longrightarrow
(U', \phi : V \to u(U'), x')
$$
is an element of $R$ which equalizes the morphisms
$(a, b, \alpha)$ and $(a', b', \alpha')$. This implies in particular
that $b = b'$. Let $d : U_2 \to U_1$ be the equalizer of $a, a'$ which
exists (see
Categories, Lemma \ref{categories-lemma-almost-finite-limits-exist}).
Moreover, $u(d) : u(U_2) \to u(U_1)$ is the equalizer of $u(a), u(a')$
hence (as $b = b'$) there is a morphism $\phi_2 : V_1 \to u(U_2)$ such that
$\phi_1 = u(d) \circ \phi_1$. Let $\delta : x_2 \to x_1$ be a strongly
cartesian morphism of $\mathcal{S}$ with $p(\delta) = u(d)$.
Now we claim that $\alpha \circ \delta = \alpha' \circ \delta$.
This is true because
$\gamma$ is strongly cartesian,
$\gamma \circ \alpha \circ \delta = \gamma \circ \alpha' \circ \delta$, and
$p(\alpha \circ \delta) = p(\alpha' \circ \delta)$.
Hence the arrow
$$
(d, \text{id}_{V_1}, \delta) :
(U_2, \phi_2 : V_1 \to u(U_2), x_2)
\longrightarrow
(U_1, \phi_1 : V_1 \to u(U_1), x_1)
$$
is an element of $R$ and equalizes $(a, b, \alpha)$ and $(a', b', \alpha')$.
Hence $R$ satisfies RMS3 as well.
\end{proof}

\end{document}
